\makeatletter
\def\rap@fontpath{../../Common/}
\makeatother
\documentclass{../../Common/Latex/Templates/rapport}
\usepackage{enumitem}

\projectname{Eclipse NMBU}
\projectlogo{../../Common/Eclipse Images/EclipseMoon}
\orglogo{../../Common/Eclipse Images/EclipseFullText}
\newcommand{\meetingline}{\textbf{Meeting:} [date], [time], [location]}

\reporttitle{Case Explanation Template}

\begin{document}

\reportmainmatter

\noindent\textbf{Document ID:} [Document ID]\\
\meetingline\\
\textbf{Companion to:} [Case number and Document ID of the Case Paper this explains]. This document has no independent legal effect: if it conflicts with the Case Paper, the Case Paper controls.

\vsection{What this Case is about}
One or two sentences stating, in plain language, what is being proposed and why the General Assembly is being asked to decide it now.

\vsection{Background}
The situation that led to the proposal: what problem it solves, what changed, or what the Association has been doing informally that this Case would formalise.

\vsection{What the proposal does}
What changes in practice if the proposal is adopted, stated as concretely as possible. Where the proposal has several distinct effects, list each as its own point:
\begin{enumerate}[label=\arabic*.]
  \item [First effect.]
  \item [Second effect.]
\end{enumerate}

\vsection{What does not change}
Anything a reader might otherwise assume is affected, but is not: existing rights, existing structures, or existing documents that remain as they are.

\vsection{Why the C-suite recommends this}
The reasoning behind the C-suite's deliberation on the Case Paper, in more depth than the Case Paper's own C-suite deliberation field allows.

\vsection{Sources and further reading}
Cross-references to the Bylaws, prior decisions, or other documents in the set that a member would need to check the proposal independently.

\end{document}
